\documentclass[a4paper,11pt]{article}
\usepackage[utf8]{inputenc}
\usepackage[czech]{babel}
\usepackage{amsmath, amssymb, amsthm}
\usepackage[margin=2cm]{geometry}
\usepackage{enumerate}

\title{\textbf{Velká písemná práce z Lineární algebry 2}}
\author{Letní semestr -- Vlastní čísla a Skalární součin}
\date{}

\begin{document}

\maketitle

\noindent \textbf{Instrukce:}
\begin{itemize}
    \item Celkový počet bodů: 100.
    \item Pro zisk zápočtu je třeba alespoň 60 bodů.
    \item \textbf{Sekce I (Ano/Ne):} K získání bodů za otázku je třeba odpovědět správně na všechny tři její podčásti.
    \item Všude, kde není řečeno jinak, značí $\langle \cdot, \cdot \rangle$ standardní skalární součin na $\mathbb{R}^n$, resp.\ $\mathbb{C}^n$.
\end{itemize}

\hrule
\vspace{0.5cm}

\section*{Část I: Ano/Ne (8 bodů)}
\textit{Rozhodněte, zda platí následující tvrzení. Není třeba zdůvodňovat. (2 body za každou správně zodpovězenou trojici)}

\begin{enumerate}[(a)]
    \item \textbf{Vlastní čísla:} Nechť $A$ je čtvercová matice řádu $n$ nad $\mathbb{R}$.
    \begin{itemize}
        \item[\textbf{ANO / NE}] Pokud $\lambda$ je vlastní číslo matice $A$ s vlastním vektorem $v$, pak $\lambda^k$ je vlastní číslo matice $A^k$ se stejným vlastním vektorem $v$ pro každé $k \in \mathbb{N}$.
        \item[\textbf{ANO / NE}] Každá matice nad $\mathbb{R}$ má alespoň jedno reálné vlastní číslo.
        \item[\textbf{ANO / NE}] Matice $A$ a $A^T$ mají stejný charakteristický polynom.
    \end{itemize}

    \item \textbf{Diagonalizovatelnost:} Nechť $A$ je čtvercová matice řádu $n$ nad $\mathbb{C}$.
    \begin{itemize}
        \item[\textbf{ANO / NE}] Pokud má $A$ právě $n$ různých vlastních čísel, pak je $A$ diagonalizovatelná.
        \item[\textbf{ANO / NE}] Pokud $A^2 = A$, pak je $A$ diagonalizovatelná.
        \item[\textbf{ANO / NE}] Pokud je algebraická násobnost každého vlastního čísla rovna jeho geometrické násobnosti, pak je $A$ diagonalizovatelná.
    \end{itemize}

    \item \textbf{Skalární součin:} Nechť $V$ je reálný vektorový prostor se skalárním součinem $\langle \cdot, \cdot \rangle$ a indukovanou normou $\|\cdot\|$.
    \begin{itemize}
        \item[\textbf{ANO / NE}] Pro libovolné $u, v \in V$ platí $|\langle u, v \rangle| \leq \|u\| \cdot \|v\|$, přičemž rovnost nastává právě tehdy, když jsou $u, v$ lineárně závislé.
        \item[\textbf{ANO / NE}] Předpis $\langle u, v \rangle := u_1 v_1 - u_2 v_2$ je skalární součin na $\mathbb{R}^2$.
        \item[\textbf{ANO / NE}] Pokud $\|u+v\|^2 = \|u\|^2 + \|v\|^2$, pak jsou vektory $u, v$ ortogonální.
    \end{itemize}

    \item \textbf{Ortogonalita a ortogonální matice:} Nechť $U \leq \mathbb{R}^n$ je podprostor a $Q \in \mathbb{R}^{n \times n}$.
    \begin{itemize}
        \item[\textbf{ANO / NE}] Platí $\dim(U) + \dim(U^\perp) = n$ a $(U^\perp)^\perp = U$.
        \item[\textbf{ANO / NE}] Matice $Q$ je ortogonální právě tehdy, když její sloupce tvoří ortonormální bázi $\mathbb{R}^n$.
        \item[\textbf{ANO / NE}] Pokud je $Q$ ortogonální, pak všechna její reálná vlastní čísla jsou rovna $\pm 1$.
    \end{itemize}
\end{enumerate}

\vspace{0.5cm}

\section*{Část II: Definice (12 bodů)}
\textit{Uveďte přesnou definici následujících pojmů (včetně předpokladů). (3 body za každou definici)}

\begin{enumerate}
    \item Vlastní číslo a vlastní vektor čtvercové matice $A$ nad tělesem $T$.
    \item Charakteristický polynom čtvercové matice. Uveďte též definici algebraické a geometrické násobnosti vlastního čísla.
    \item Skalární součin na vektorovém prostoru $V$ nad $\mathbb{R}$ (resp.\ $\mathbb{C}$).
    \item Ortogonální doplněk $U^\perp$ podmnožiny $U$ prostoru $V$ se skalárním součinem.
\end{enumerate}

\vspace{0.5cm}

\section*{Část III: Jednoduché příklady (15 bodů)}
\textit{Vypočtěte nebo určete. Stačí uvést správný výsledek. (3 body za každý příklad)}

\begin{enumerate}
    \item Najděte všechna (i komplexní) vlastní čísla matice 
    $$A = \begin{pmatrix} 2 & -1 \\ 5 & -2 \end{pmatrix} \text{ nad } \mathbb{R}.$$
    
    \item Rozhodněte, zda je matice $\begin{pmatrix} 3 & 1 \\ 0 & 3 \end{pmatrix}$ diagonalizovatelná nad $\mathbb{R}$. Odpověď stručně zdůvodněte.
    
    \item V prostoru $\mathbb{R}^4$ se standardním skalárním součinem vypočtěte kosinus úhlu mezi vektory $u = (1, 2, 2, 0)^T$ a $v = (0, 1, 1, 1)^T$.
    
    \item Gram--Schmidtovým ortogonalizačním procesem ortogonalizujte dvojici vektorů $v_1 = (1, 1, 0)^T$, $v_2 = (1, 0, 1)^T$ v $\mathbb{R}^3$ se standardním skalárním součinem.
    
    \item Najděte ortogonální doplněk podprostoru $U = \mathrm{span}\{(1,2,1,0)^T, (0,1,1,1)^T\}$ v $\mathbb{R}^4$.
\end{enumerate}

\vspace{0.5cm}

\section*{Část IV: Formulace tvrzení (12 bodů)}
\textit{Zformulujte přesně následující věty/tvrzení (není třeba dokazovat). (3 body za každé tvrzení)}

\begin{enumerate}
    \item Cayleyova--Hamiltonova věta.
    \item Věta o nutné a postačující podmínce diagonalizovatelnosti matice (pomocí násobností vlastních čísel).
    \item Cauchyova--Schwarzova nerovnost (včetně podmínky pro rovnost).
    \item Věta o existenci ortonormální báze v konečněrozměrném prostoru se skalárním součinem (Gram--Schmidtova věta).
\end{enumerate}

\newpage

\section*{Část V: Početní příklady s postupem (15 bodů)}
\textit{Vyřešte následující úlohy. Uveďte podrobný postup řešení. (5 bodů za každý příklad)}

\begin{enumerate}
    \item Najděte všechna vlastní čísla a odpovídající vlastní vektory matice 
    $$A = \begin{pmatrix} 1 & 2 & 2 \\ 2 & 1 & 2 \\ 2 & 2 & 1 \end{pmatrix}$$
    nad $\mathbb{R}$. Rozhodněte, zda je $A$ diagonalizovatelná, a pokud ano, najděte regulární matici $P$ a diagonální matici $D$ takové, že $P^{-1} A P = D$.
    
    \item V prostoru $\mathbb{R}^4$ se standardním skalárním součinem proveďte Gram--Schmidtovu ortonormalizaci vektorů 
    $$v_1 = (1, 1, 0, 0)^T, \quad v_2 = (1, 0, 1, 0)^T, \quad v_3 = (1, 0, 0, 1)^T.$$
    
    \item Nechť $U \leq \mathbb{R}^4$ je podprostor s bází $\{(1, 1, 0, 0)^T, (0, 0, 1, 1)^T\}$ (se standardním skalárním součinem). Najděte ortogonální projekci vektoru $v = (1, 2, 3, 4)^T$ do podprostoru $U$ a spočtěte vzdálenost vektoru $v$ od podprostoru $U$.
\end{enumerate}

\vspace{0.5cm}

\section*{Část VI: Důkazy jednodušších tvrzení (20 bodů)}
\textit{Dokažte následující tvrzení. (5 bodů za každý důkaz)}

\begin{enumerate}
    \item Nechť $A$ je čtvercová matice nad tělesem $T$ a $\lambda_1, \ldots, \lambda_k$ jsou navzájem různá vlastní čísla $A$ s odpovídajícími vlastními vektory $v_1, \ldots, v_k$. Dokažte, že vektory $v_1, \ldots, v_k$ jsou lineárně nezávislé.
    
    \item Dokažte Cauchyovu--Schwarzovu nerovnost: pro libovolné vektory $u, v$ v reálném prostoru se skalárním součinem platí $|\langle u, v \rangle| \leq \|u\| \cdot \|v\|$.
    
    \item Nechť $U$ je konečněrozměrný podprostor prostoru $V$ se skalárním součinem. Dokažte, že $V = U \oplus U^\perp$ (tj.\ že $U \cap U^\perp = \{o\}$ a $V = U + U^\perp$).
    
    \item Dokažte, že pro každou ortogonální matici $Q \in \mathbb{R}^{n \times n}$ platí $|\det(Q)| = 1$. Ukažte též, že podobné matice mají stejný charakteristický polynom.
\end{enumerate}

\vspace{0.5cm}

\section*{Část VII: Úlohy na zamyšlení (18 bodů)}
\textit{Vyřešte nebo dokažte následující náročnější úlohy. (6 bodů za každou úlohu)}

\begin{enumerate}
    \item \textbf{Vlastní čísla symetrických matic:} Nechť $A \in \mathbb{R}^{n \times n}$ je symetrická matice (tj.\ $A^T = A$). Dokažte, že:
    \begin{enumerate}[(a)]
        \item všechna vlastní čísla matice $A$ jsou reálná,
        \item vlastní vektory odpovídající různým vlastním číslům jsou na sebe kolmé (vzhledem ke standardnímu skalárnímu součinu).
    \end{enumerate}
    
    \item \textbf{Besselova nerovnost:} Nechť $V$ je prostor se skalárním součinem a $\{e_1, e_2, \ldots, e_k\}$ je ortonormální systém vektorů ve $V$. Dokažte, že pro každý vektor $v \in V$ platí
    $$
    \sum_{i=1}^{k} |\langle v, e_i \rangle|^2 \leq \|v\|^2.
    $$
    Určete, kdy nastává rovnost.
    
    \item \textbf{Stopa a vlastní čísla:} Nechť $A \in \mathbb{C}^{n \times n}$ a $\lambda_1, \ldots, \lambda_n$ jsou vlastní čísla matice $A$ (počítaná s algebraickými násobnostmi). Dokažte, že
    $$
    \mathrm{tr}(A) = \sum_{i=1}^{n} \lambda_i \quad \text{a} \quad \det(A) = \prod_{i=1}^{n} \lambda_i.
    $$
    \textit{Nápověda: použijte vztah charakteristického polynomu a jeho koeficientů.}
\end{enumerate}

\end{document}